\documentclass[11pt,a4paper]{article}
\usepackage[utf8]{inputenc}
\usepackage[T1]{fontenc}
\usepackage[spanish]{babel}
\usepackage{hyperref}
\usepackage{geometry}
\geometry{margin=2.5cm}
\usepackage{parskip}

\title{012 la limitacion declarada}
\author{Alejandro Valenzuela Rivera\\\small Universidad Autónoma de Nuevo León}
\date{\today}
\begin{document}
\maketitle

\begin{quote}
\textbf{Sesgo declarado.} Este trabajo se sostiene desde el sesgo declarado de su autor, y no finge neutralidad. Lo que afirma se puede tumbar, y ese es el punto: el falsificador está escrito, no escondido.
\end{quote}

\section*{Resumen}
Un ente que no sabe lo que no sabe contamina.

\section{Por qué se sostiene}
Saber las limitaciones es la unica defensa contra la memoria absoluta.

\section{De dónde viene}
La autoridad, y el repo: el borde que el ente declara.

\section{El caso}
Ninguna: es la regla del borde declarado.

\section{Método}
Un ente que dice que no sabe, antes de que se le pregunte.

\section{Límites, y qué lo tumba}
Un ente que se presenta completo y contesta como si supiera.

\section{Lo que queda abierto}
\textquestiondown{}Que pasa cuando el ente llena lo que no sabe?

\section*{Cómo se leyó esta tesis}
Esto se escribió en el repositorio man-made-empathy, del que este candado es una pieza: candados/04 el lenguaje/012 la limitacion declarada.md. La máquina que lo compuso lo hizo desde este mismo archivo, y no agregó nada que no estuviera escrito. Si algo de esto es falso, se dice en el mismo sitio, y se corrige con fecha.

\end{document}
